% ============================================================================
%  Growth & Remodeling seminar — shared beamer preamble
%  Adapted from the YaleCBL-Teaching/BENG_5570 slide build system.
%  Two build modes are selected by \teachertrue / \teacherfalse in the
%  master files (slides-student.tex / slides-teacher.tex):
%    student  -> key equations shown as EMPTY numbered boxes to fill on iPad
%    teacher  -> the same boxes are filled with the actual equation
% ============================================================================
\usepackage[utf8]{inputenc}
\usepackage[T1]{fontenc}
\usepackage{amsmath,amssymb,amsfonts,mathtools}
\usepackage{bm}
\usepackage{graphicx}
\usepackage{booktabs}
\usepackage{array}
\usepackage{tikz}
\usetikzlibrary{positioning,arrows.meta,calc}
\usepackage[export]{adjustbox}
\usepackage{xcolor}

% ---- sans-serif math+text (matches BENG_5570 house style) ------------------
\usepackage{helvet}
\renewcommand{\familydefault}{\sfdefault}
\usefonttheme{professionalfonts}

% ---- Yale-blue palette -----------------------------------------------------
\definecolor{YaleBlue}{RGB}{0,53,107}
\definecolor{YaleLightBlue}{RGB}{40,109,192}
\definecolor{AccentRed}{RGB}{171,38,38}
\definecolor{BoxGray}{RGB}{120,120,120}
\definecolor{SlideFrame}{RGB}{0,53,107}

% ---- beamer theme ----------------------------------------------------------
\usetheme{default}
\usecolortheme{default}
\setbeamercolor{frametitle}{fg=YaleBlue}
\setbeamercolor{title}{fg=YaleBlue}
\setbeamercolor{structure}{fg=YaleLightBlue}
\setbeamercolor{item}{fg=YaleBlue}
\setbeamertemplate{navigation symbols}{}
\setbeamertemplate{itemize items}[circle]
\setbeamerfont{frametitle}{series=\bfseries}
\setbeamersize{text margin left=1.2em,text margin right=1.2em}

% footline: course tag + slide number
% Lower-left carries the slide's reference (set via \srccite), not a running
% title; lower-right the frame number. The stored reference resets each shipout.
\setbeamertemplate{footline}{%
  \leavevmode\hbox{%
    \begin{beamercolorbox}[wd=\paperwidth,dp=1.2ex,leftskip=1.2em,rightskip=1.2em]{}%
      {\scriptsize\color{BoxGray}%
        \begin{minipage}[b]{\dimexpr\paperwidth-5.2em\relax}\itshape\sliderefstore\end{minipage}%
        \hfill
        \begin{minipage}[b]{1.6em}\raggedleft\upshape\insertframenumber\end{minipage}}%
    \end{beamercolorbox}}%
  \gdef\sliderefstore{}%
}

% ============================================================================
%  Fill-in-the-blank equation mechanism
% ============================================================================
% \teacher toggle set in the master file.
\newif\ifteacher

% Height register for a student blank box.
% Usage in slides:
%   \slboxeq{<box height>}{<eq label>}{<equation body>}
%     student : empty framed box of the given height, right-numbered (label)
%     teacher : the equation, boxed and numbered (same label)
% The label is shared, so \eqref works identically in both builds.
\newcommand{\slboxeq}[3]{%
  \ifteacher
    % teacher: show the boxed, numbered equation
    \begin{equation}\label{#2}
      \boxed{\displaystyle #3}
    \end{equation}
  \else
    % student: reserve an empty framed box carrying the equation number
    \begin{equation}\label{#2}\end{equation}%
    \vspace{-2.2\baselineskip}%
    \begin{center}%
    \setlength{\fboxrule}{0.8pt}\setlength{\fboxsep}{4pt}%
    \textcolor{SlideFrame}{\framebox[0.92\linewidth][c]{%
      \rule{0pt}{#1}\rule{0pt}{0pt}%
    }}%
    \end{center}%
    \vspace{-0.3\baselineskip}%
  \fi
}

% Inline "given" equation (always shown in both builds, unnumbered box-free).
\newcommand{\giveneq}[1]{\begin{equation*}#1\end{equation*}}

% ============================================================================
%  3D <-> 0D presentation convention
%  Every core relation is shown first in its general 3D (tensor) form, then in
%  the 0D scalar reduction actually used in this course. A small colored badge
%  marks which regime a line belongs to; \formlab places the badge + a short
%  caption on a tight line just above the equation it introduces.
% ============================================================================
\definecolor{TagD}{RGB}{62,78,138}    % indigo/slate  -> 3D (general tensor)
\definecolor{TagZ}{RGB}{17,122,116}   % teal          -> 0D (scalar reduction)
\newcommand{\dpill}[2]{\tikz[baseline]{%
  \node[fill=#1,text=white,rounded corners=2.5pt,inner xsep=4pt,inner ysep=1.6pt,
        font=\bfseries\scriptsize,anchor=base]{#2};}}
\newcommand{\tagd}{\dpill{TagD}{3D}}
\newcommand{\tagz}{\dpill{TagZ}{0D}}
% \formlab{<badge>}{<short caption>} — tight labeled line before an equation
\newcommand{\formlab}[2]{\par\smallskip\noindent
  {#1}\;{\footnotesize\color{black!62}#2}\par\nobreak\vspace{-0.2em}}
% convenience wrappers
\newcommand{\dlab}[1]{\formlab{\tagd}{#1}}   % 3D general (tensor) form
\newcommand{\zlab}[1]{\formlab{\tagz}{#1}}   % 0D scalar reduction

% ---- per-slide source citation, required format ----------------------------
% "First, ..., Last, short journal (year)"
% The reference is not printed in the slide body; it is stashed and shown in the
% lower-left footline (see the footline template above). Reset each shipout.
\newcommand{\sliderefstore}{}
\newcommand{\srccite}[1]{\gdef\sliderefstore{#1}}
% bottom-right floating credit for a figure
\newcommand{\figcredit}[1]{%
  \begin{flushright}\vspace{-0.6em}{\tiny\color{BoxGray} #1}\end{flushright}%
}

% ---- convenience figure command --------------------------------------------
\newcommand{\slidefig}[2][0.8\linewidth]{%
  \begin{center}\includegraphics[width=#1,height=0.62\textheight,keepaspectratio]{figures/#2}\end{center}%
}

% ---- big-figure slide: one large centred figure + short caption ------------
% \bigfig{file}{one-line caption}  — figure dominates the slide, minimal text.
% Captions under figures are intentionally suppressed (#2/#3 ignored).
\newcommand{\bigfig}[2]{%
  \vfill
  \begin{center}%
    \includegraphics[width=0.92\linewidth,height=0.66\textheight,keepaspectratio]{figures/#1}%
  \end{center}%
  \vfill
}
% variant with adjustable width for tall figures
\newcommand{\bigfigw}[3]{%
  \vfill
  \begin{center}%
    \includegraphics[width=#2,height=0.66\textheight,keepaspectratio]{figures/#1}%
  \end{center}%
  \vfill
}

% ---- math shorthands used across the deck ----------------------------------
\newcommand{\dd}{\mathrm{d}}
\newcommand{\Ftens}{\bm{F}}
\newcommand{\sig}{\sigma}
\newcommand{\sigbar}{\bar{\sigma}}
\newcommand{\sigh}{\bar{\sigma}_h}
\newcommand{\lam}{\lambda}
\newcommand{\Ups}{\Upsilon}
\newcommand{\kd}{k_d}

% ============================================================================
%  Interactive-exercise format
%  A standard "Your turn" question frame (Question / Resources / Hints) and a
%  matching worked "Answer" frame. Used throughout so there is always something
%  to do every few slides.
% ============================================================================
\usepackage{tcolorbox}
\usepackage{fontawesome5}
\tcbuselibrary{skins,breakable}
\definecolor{QBlue}{RGB}{40,109,192}
\definecolor{ABlue}{RGB}{22,120,70}   % green accent for answers
\definecolor{QFill}{RGB}{236,242,250}
\definecolor{AFill}{RGB}{234,246,238}

% \begin{yourturn}{<short title>} ... \fields ... \end{yourturn}
\newtcolorbox{yourturnbox}[1]{%
  enhanced, breakable, colback=QFill, colframe=QBlue, boxrule=1pt,
  arc=2pt, left=6pt,right=6pt,top=4pt,bottom=4pt,
  fonttitle=\bfseries\color{white},
  title={\faPencilRuler~~Your turn: #1}}
\newtcolorbox{answerbox}[1]{%
  enhanced, breakable, colback=AFill, colframe=ABlue, boxrule=1pt,
  arc=2pt, left=6pt,right=6pt,top=4pt,bottom=4pt,
  fonttitle=\bfseries\color{white},
  title={\faCheckCircle~~Answer: #1}}

% field labels inside a yourturn box
\newcommand{\qq}[1]{\par\smallskip{\color{QBlue}\bfseries Question.}\ #1}
\newcommand{\qres}[1]{\par\smallskip{\color{QBlue}\bfseries Resources.}\ #1}
\newcommand{\qhint}[1]{\par\smallskip{\color{QBlue}\bfseries Hints.}\ #1}
\newcommand{\qkind}[1]{}


% ===========================================================================
%  Idealized blood vessel drawn in TikZ (thick-walled cylinder cross-section)
%    \idealvessel        -> full annotated version (P, Q, r, h, sigma_theta, tau_w)
% ============================================================================
\newcommand{\idealvessel}[1][1.0]{%
\begin{tikzpicture}[scale=#1]
  % wall (annulus) and lumen
  \fill[YaleLightBlue!18] (0,0) circle (2.55);
  \fill[white] (0,0) circle (1.8);
  \draw[YaleBlue,very thick] (0,0) circle (2.55);
  \draw[YaleBlue,very thick] (0,0) circle (1.8);
  % inner radius r
  \draw[-{Latex[length=2mm]},thick] (0,0) -- (128:1.8) node[midway,above,sloped]{\small $r$};
  \fill (0,0) circle (1.2pt);
  % wall thickness h
  \draw[{Latex[length=1.6mm]}-{Latex[length=1.6mm]},thick,AccentRed]
     (-18:1.8) -- (-18:2.55);
  \node[AccentRed] at (-18:2.9) {\small $h$};
  % lumen pressure P (radial arrows outward on inner wall)
  \foreach \a in {30,70,110,150,190,230,270,310}{
     \draw[-{Latex[length=1.6mm]},YaleLightBlue] (\a:0.55) -- (\a:1.55);}
  \node[YaleLightBlue] at (0,0.0) {\small $P$};
  % blood flow Q (into the page / axial) + WSS tau_w at inner wall
  \node[YaleLightBlue] at (0,-0.75) {\small $Q$};
  \draw[-{Latex[length=1.6mm]},AccentRed] (60:1.75) arc (60:20:1.75);
  \node[AccentRed] at (40:1.45) {\scriptsize $\tau_w$};
  % circumferential stress sigma_theta (tangential arrows in the wall)
  \draw[-{Latex[length=1.6mm]},YaleBlue] (95:2.18) arc (95:60:2.18);
  \draw[-{Latex[length=1.6mm]},YaleBlue] (85:2.18) arc (85:120:2.18);
  \node[YaleBlue] at (90:2.02) {\scriptsize $\sigma_\theta$};
\end{tikzpicture}}

% ============================================================================
%  Reference -> current configuration mapping (TikZ)  \configmap
% ============================================================================
% ============================================================================
%  Configuration-mapping diagrams (TikZ, grayscale)
%    \configmap  reference -> current (deformation gradient F)
%    \kinmap     kinematic growth multiplicative split  F = F_e F_g
%    \cmmmap     constrained mixture: cohorts, each its own natural config
%  Shared style: stress-free/reference bodies light gray, loaded/current
%  bodies darker gray; labels placed clear of the blobs to avoid overlap.
% ============================================================================
\tikzset{
  refbody/.style={draw=black!70,very thick,fill=black!8},
  intbody/.style={draw=black!70,very thick,fill=black!16},
  curbody/.style={draw=black!80,very thick,fill=black!26},
  cfgmaparrow/.style={-{Latex[length=2.2mm]},thick,black!75},
  cfgbase/.style={-{Latex[length=1.2mm]},black!60},
}

% ---- reference -> current --------------------------------------------------
\newcommand{\configmap}[1][1.0]{%
\begin{tikzpicture}[scale=#1,every node/.style={font=\small}]
  % reference body (left)
  \draw[refbody] plot[smooth cycle,tension=0.9]
    coordinates {(0,0)(1.0,0.55)(1.8,0.15)(1.55,-0.85)(0.55,-1.05)(-0.15,-0.5)};
  \fill (0.85,-0.2) circle (1.4pt);
  \node[anchor=south west,font=\footnotesize] at (0.9,-0.2) {$\bm{X}$};
  \draw[cfgbase] (0.85,-0.2)--(1.2,-0.2);
  \draw[cfgbase] (0.85,-0.2)--(0.85,0.15);
  \node[font=\footnotesize] at (0.8,-1.5) {reference $\mathcal{B}_0$};
  % current body (right)
  \begin{scope}[shift={(5.0,-0.05)}]
    \draw[curbody] plot[smooth cycle,tension=0.9]
      coordinates {(0,0.2)(1.25,0.7)(2.25,0.2)(2.1,-1.0)(1.0,-1.45)(-0.1,-0.75)};
    \fill (1.25,-0.3) circle (1.4pt);
    \node[anchor=south west,font=\footnotesize] at (1.3,-0.3) {$\bm{x}$};
    \draw[cfgbase] (1.25,-0.3)--(1.65,-0.18);
    \draw[cfgbase] (1.25,-0.3)--(1.32,0.12);
    \node[font=\footnotesize] at (1.05,-1.9) {current $\mathcal{B}_t$};
  \end{scope}
  % motion + line-element mapping
  \draw[cfgmaparrow] (2.0,0.55) to[bend left=20]
    node[above,font=\small]{$\bm{\varphi}(t,\bm{X})$} (5.1,0.75);
  \draw[cfgmaparrow,black!55] (2.0,-0.25) to[bend right=10]
    node[below,font=\footnotesize]{$\dd\bm{x}=\bm{F}\,\dd\bm{X}$} (6.0,-0.45);
\end{tikzpicture}}

% ---- kinematic growth: F = F_e F_g, with INCOMPATIBLE grown configuration ---
% Reference body is cut into pieces; each piece grows stress-free on its own, so
% the grown pieces no longer tile (overlaps + gaps = incompatible). The elastic
% step F_e forces them back into one connected body (residual stress).
\newcommand{\kinmap}[1][1.0]{%
\begin{tikzpicture}[scale=#1,every node/.style={font=\small},
    piece/.style={draw=black!70,thick,fill=black!10},
    gpiece/.style={draw=black!70,thick,dashed,fill=black!16},
    cpiece/.style={draw=black!80,thick,fill=black!24}]
  % --- reference B_0: 2x2 compatible tiling (left) ---
  \begin{scope}[shift={(0,0)}]
    \foreach \x in {0,0.9} \foreach \y in {-0.9,0}{
      \draw[piece] (\x,\y) rectangle ++(0.9,0.9);}
    \node[font=\footnotesize] at (0.9,-1.35) {reference $\mathcal{B}_0$};
  \end{scope}
  % --- grown, stress-free, INCOMPATIBLE (middle): each piece enlarged about
  %     its own centre + jittered so they overlap and leave gaps ---
  \begin{scope}[shift={(3.7,0.05)}]
    \draw[gpiece] (-0.12,0.02) rectangle ++(1.14,1.14);   % top-left, shifted
    \draw[gpiece] (1.02,0.12)  rectangle ++(1.14,1.14);   % top-right, overlaps
    \draw[gpiece] (0.06,-1.18) rectangle ++(1.14,1.14);   % bot-left
    \draw[gpiece] (1.16,-1.08) rectangle ++(1.14,1.14);   % bot-right, gap
    \node[font=\footnotesize,align=center] at (1.05,-1.7)
      {grown, stress-free\\\emph{incompatible} $\mathcal{B}_g$};
  \end{scope}
  % --- current, loaded, COMPATIBLE (right): reassembled connected body ---
  \begin{scope}[shift={(7.6,0)}]
    \foreach \x in {0,1.02} \foreach \y in {-1.02,0}{
      \draw[cpiece] (\x,\y) rectangle ++(1.02,1.02);}
    \node[font=\footnotesize] at (1.02,-1.42) {current $\mathcal{B}_t$};
  \end{scope}
  % --- arrows ---
  \draw[cfgmaparrow] (1.95,0.15) to[bend left=16]
    node[above,font=\small]{$\bm{F}_g$ (growth, prescribed)} (3.7,0.35);
  \draw[cfgmaparrow] (6.4,1.35) to[bend left=12]
    node[above,font=\small,pos=0.4]{$\bm{F}_e$ (elastic, restores fit)} (7.7,0.7);
  \draw[cfgmaparrow,black!55] (0.5,-2.35) to[bend right=8]
    node[below,font=\footnotesize,pos=0.5]{$\bm{F}=\bm{F}_e\,\bm{F}_g$} (8.6,-2.15);
\end{tikzpicture}}

% ---- Humphrey (2021) Fig. 1: finite deformations in a constrained mixture -----
% Horizontal barrel of stacked ellipse cross-sections = tissue evolving from
% kappa(0) -> kappa(tau) -> kappa(s); a deposited constituent sits at its own
% stress-free natural config kappa_n^alpha(tau) above kappa(tau), carried in by
% the deposition stretch G^alpha(tau); F_{n(tau)}^alpha(s) is the constituent-
% specific deformation; F(tau), F(s) are the tissue-level deformations.
\newcommand{\humphreygr}[1][1.0]{%
\begin{tikzpicture}[scale=#1,every node/.style={font=\small},
    ell/.style={draw=black!65,thick,fill=black!12},
    ellface/.style={draw=black!75,very thick,fill=black!18},
    grarr/.style={-{Latex[length=2.4mm]},very thick,black!78},
    tisarr/.style={-{Latex[length=2.2mm]},thick,black!70}]
  \def\ry{0.85}   % ellipse vertical radius
  \def\rx{0.30}   % ellipse horizontal radius
  % --- barrel: many faint ellipses stacked left->right ---
  \foreach \i in {0,1,...,16}{
    \pgfmathsetmacro\xx{0.6+\i*0.32}
    \draw[ell,fill opacity=0.85] (\xx,0) ellipse [x radius=\rx, y radius=\ry];}
  % --- three labeled cross-sections: kappa(0), kappa(tau), kappa(s) ---
  \draw[ellface] (0.6,0)  ellipse [x radius=\rx, y radius=\ry];   % kappa(0)
  \draw[ellface] (3.16,0) ellipse [x radius=\rx, y radius=\ry];   % kappa(tau)  (i=8)
  \draw[ellface] (5.72,0) ellipse [x radius=\rx, y radius=\ry];   % kappa(s)    (i=16)
  \node[font=\footnotesize] at (0.6,0)  {$\kappa(0)$};
  \node[font=\footnotesize] at (3.16,0) {$\kappa(\tau)$};
  \node[font=\footnotesize] at (5.72,0) {$\kappa(s)$};
  % dash-dot centre axis extending both ends
  \draw[densely dashdotted,black!55] (-0.5,0)--(0.28,0);
  \draw[densely dashdotted,black!55] (6.04,0)--(6.9,0);
  \node[align=center,font=\footnotesize,black!75] at (7.0,-0.55)
    {Growth \&\\Remodeling};
  % --- deposited constituent at its stress-free natural config (top) ---
  \node[circle,draw=black!70,thick,fill=black!22,minimum size=6pt,inner sep=0pt]
    (kn) at (2.55,2.25) {};
  \node[anchor=east,font=\footnotesize] at (2.42,2.35) {$\kappa_n^{\alpha}(\tau)$};
  % two small chained blobs suggesting the constituent chain
  \fill[black!22,draw=black!70] (2.68,1.95) circle (2.2pt);
  \fill[black!22,draw=black!70] (2.82,1.68) circle (2.0pt);
  % deposition stretch G^alpha(tau): from natural config down into kappa(tau)
  \draw[grarr] (2.78,1.55) to[bend left=25]
    node[left,font=\small,pos=0.55]{$\bm{G}^{\alpha}(\tau)$} (3.02,0.7);
  % constituent-specific deformation F_{n(tau)}^alpha(s): natural config -> right
  \draw[grarr] (2.95,2.35) to[bend left=20]
    node[above,font=\small,pos=0.5]{$\bm{F}^{\alpha}_{n(\tau)}(s)$} (5.5,1.15);
  % --- tissue-level deformations underneath ---
  \draw[tisarr] (0.7,-1.05) to[bend right=22]
    node[below,font=\small,pos=0.5]{$\bm{F}(\tau)$} (3.06,-1.05);
  \draw[tisarr] (0.7,-1.55) to[bend right=20]
    node[below,font=\small,pos=0.55]{$\bm{F}(s)$} (5.62,-1.55);
\end{tikzpicture}}

% ---- Cyron (2016) Fig. 2 (right): mechanical analog for mass turnover ---------
% A motor element (circle) exerting the prestress sigma_pre in PARALLEL with a
% Maxwell material = dashpot (time constant rho-dot0/rho0 + 1/T) in SERIES with a
% (nonlinear) spring of stiffness C. Load sigma applied top and bottom.
\newcommand{\maxwellelement}[1][1.0]{%
\begin{tikzpicture}[scale=#1,line width=1pt,black!85,every node/.style={font=\small}]
  % geometry
  \def\xm{0}      % Maxwell arm x
  \def\xp{3.2}    % motor arm x
  \def\ytop{4.6}
  \def\ybot{0}
  % top & bottom load leads with sigma
  \draw (\xm,\ytop)--(\xp,\ytop);                 % top rail
  \draw (\xm,\ybot)--(\xp,\ybot);                 % bottom rail
  \coordinate (topmid) at ($(\xm,\ytop)!0.5!(\xp,\ytop)$);
  \coordinate (botmid) at ($(\xm,\ybot)!0.5!(\xp,\ybot)$);
  \draw (topmid)--++(0,0.7) node[above,font=\normalsize]{$\bm{\sigma}^i$};
  \draw[-{Latex[length=2mm]}] (topmid)++(0,0.62)--++(0,0.28);
  \draw (botmid)--++(0,-0.7) node[below,font=\normalsize]{$\bm{\sigma}^i$};
  \draw[-{Latex[length=2mm]}] (botmid)++(0,-0.62)--++(0,-0.28);
  % ===== Maxwell arm (left): dashpot on top, spring below =====
  % lead from top rail down to dashpot
  \draw (\xm,\ytop)--(\xm,4.15);
  % dashpot: cylinder + piston
  \draw (\xm-0.32,4.15)--(\xm-0.32,3.5)--(\xm+0.32,3.5)--(\xm+0.32,4.15); % cylinder (open top)
  \draw[line width=3pt] (\xm-0.22,3.95)--(\xm+0.22,3.95);                 % piston plate
  \draw (\xm,4.15)--(\xm,3.95);                                          % piston rod
  \node[right,font=\footnotesize,align=left] at (\xm+0.5,3.82)
    {$\dfrac{\dot\varrho_0^{\,i}}{\varrho_0^{\,i}}+\dfrac{1}{T^{i}}$};
  % series link dashpot -> spring
  \draw (\xm,3.5)--(\xm,3.1);
  % spring (zig-zag) 3.1 -> 1.5
  \draw (\xm,3.1)
    --(\xm-0.28,2.95)--(\xm+0.28,2.75)--(\xm-0.28,2.55)--(\xm+0.28,2.35)
    --(\xm-0.28,2.15)--(\xm+0.28,1.95)--(\xm-0.28,1.75)--(\xm,1.6)--(\xm,1.5);
  \node[right,font=\footnotesize] at (\xm+0.42,2.35) {$\mathbb{C}^{i}$};
  % lead spring -> bottom rail
  \draw (\xm,1.5)--(\xm,\ybot);
  % ===== Motor / prestress arm (right): circle with sigma_pre =====
  \draw (\xp,\ytop)--(\xp,3.15);
  \draw (\xp,3.15) circle (0.42);
  \node[left,font=\footnotesize] at (\xp-0.6,3.15) {$\bar{\bm{\sigma}}^{i}_{\text{pre}}$};
  \draw (\xp,2.73)--(\xp,\ybot);
\end{tikzpicture}}

% ---- constrained mixture: cohorts, each its own natural configuration ------
\newcommand{\cmmmap}[1][1.0]{%
\begin{tikzpicture}[scale=#1,every node/.style={font=\small}]
  % three small stress-free cohort configs (left, stacked)
  \foreach \y/\lab in {1.7/{\tau=0}, 0.0/{\tau}, -1.7/{s}} {
    \begin{scope}[shift={(0,\y)}]
      \draw[refbody] plot[smooth cycle,tension=0.85]
        coordinates {(0,0)(0.55,0.32)(0.95,0.05)(0.82,-0.5)(0.32,-0.62)(-0.06,-0.3)};
      \node[font=\scriptsize,anchor=east] at (-0.15,-0.15) {$\lab$};
    \end{scope}}
  \node[font=\footnotesize,align=center] at (0.45,-2.75)
    {cohorts, each at\\deposition stretch $G^k$};
  % shared current mixture (right)
  \begin{scope}[shift={(5.0,0)}]
    \draw[curbody] plot[smooth cycle,tension=0.9]
      coordinates {(0,0.4)(1.3,1.0)(2.3,0.35)(2.15,-1.1)(1.0,-1.6)(-0.1,-0.85)};
    \node[font=\footnotesize] at (1.05,-2.1) {current mixture $\mathcal{B}_t$};
    \node[font=\scriptsize,align=center,black!70] at (1.05,-0.35)
      {constrained to\\deform together};
  \end{scope}
  % arrows from each cohort into the current mixture
  \draw[cfgmaparrow] (1.0,1.55) to[bend left=10]
    node[above,font=\footnotesize]{$\bm{F}_{n(\tau)}$} (5.0,0.5);
  \draw[cfgmaparrow] (1.0,-0.15) -- (4.9,-0.25);
  \draw[cfgmaparrow] (1.0,-1.85) to[bend right=10] (5.0,-1.15);
\end{tikzpicture}}


% ---- compact coding-exercise field (adds a monospace "run this" line) ------
\newcommand{\qrun}[1]{\par\smallskip{\color{QBlue}\bfseries Run.}\ \texttt{#1}}
